\chapter{複素平面と漸化式}\label{章：複素平面と漸化式}
% 終章の最終ページだけは左右の段を均等に終える。
\begin{multicols}{2}
実数の数列に複素数を使うと,\,一つの値の背後にある平面上の動きが見える.
まず積と商を図形として読み,\,その後で1次分数型と3項間漸化式へ進もう.
\section{極形式と複素平面}\label{節：極形式と複素平面}
\subsection{複素数を点として表す}

$i^2=-1$を満たす虚数単位$i$を用い,\,複素数を$z=x+iy$と書く.
$x,y$は実数である.
$x$を実部,$y$を虚部といい,\,平面上の点$(x,y)$と対応させる.
横軸を実軸,\,縦軸を虚軸と呼ぶこの平面が\textbf{複素平面}である.

複素数の和は成分ごとの和だから,\,平面の矢印の和と一致する.
積は$i^2=-1$を使って計算する.
例えば
\[
 i(x+iy)=-y+ix
\]
であり,\,点$(x,y)$は$(-y,x)$へ移る.
これは原点のまわりの反時計回りの直角回転である.

$z=x+iy$に対して$\overline z=x-iy$を\textbf{共役複素数}という.
実軸に関する鏡映に当たり,
\[
 z\overline z=x^2+y^2=|z|^2
\]
となる. $|z|=\sqrt{x^2+y^2}$を絶対値という.
$z\ne0$なら
\[
 \frac1z=\frac{\overline z}{|z|^2}
\]
だから,\,複素数でも零以外で割ることができる.

\subsection{極形式と積}

零でない$z$を,\,原点からの距離$\rho>0$と角度$\theta$で表す.
\[
 z=\rho(\cos\theta+i\sin\theta).
\]
これを\textbf{極形式}という.
角度$\theta$は偏角と呼ばれ,$2\pi$の整数倍を加えても同じ点を表す.
零には偏角を定めない.

加法公式
\[
 \begin{aligned}
 \cos(\theta+\phi)&=\cos\theta\cos\phi-\sin\theta\sin\phi,\\
 \sin(\theta+\phi)&=\sin\theta\cos\phi+\cos\theta\sin\phi
 \end{aligned}
\]
を使うと,
\[
 \begin{aligned}
 &(\cos\theta+i\sin\theta)(\cos\phi+i\sin\phi)\\
 &\qquad=\cos(\theta+\phi)+i\sin(\theta+\phi).
 \end{aligned}
\]
従って積では距離を掛け,\,角度を足す.
商では距離を割り,\,角度を引く.

$\lambda=\rho(\cos\theta+i\sin\theta)$を固定して掛けると,\,原点からの距離は$\rho$倍,\,角度は$\theta$だけ進む.
$\rho=1$なら一定角の回転,$\rho<1$なら縮小を伴う回転となる.
\NewFigRotation

積の公式を繰り返せば,\,非負整数$m$について
\[
 (\cos\theta+i\sin\theta)^m
 =\cos(m\theta)+i\sin(m\theta).
\]
これが\textbf{ド・モアブルの定理}である.
本書では$\cos\theta+i\sin\theta$を$e^{i\theta}$とも略記する.
ここでは回転を表す記号として使い,\,複素指数関数の一般論は必要としない.

\subsection{商は長さの比と角度の差}

二点$\alpha,\beta$を基準にして,\,点$z$の位置を
\[
 w=\frac{z-\alpha}{z-\beta}\qquad(z\ne\beta)
\]
で記録してみよう.
絶対値は
\[
 |w|=\frac{|z-\alpha|}{|z-\beta|}
\]
だから,\,二つの基準点からの距離の比を表す.
偏角は,\,ベクトル$z-\beta$から$z-\alpha$への有向角である.
二つの矢印を同じ始点に平行移動して角度を比較してもよい.
\NewFigRatio

この比が一定倍率で変わるなら,\,距離の比と角度の差の変化を同時に追える.
次節で,1次分数型に現れる置き換えがこの形になる.

\begin{bookpoint}
複素数の積は\textbf{倍率と回転角}をまとめて記録する.
複素数の商は\textbf{長さの比と角度の差}をまとめて記録する.
\end{bookpoint}

\section{1次分数変換}\label{節：一次分数変換}
\subsection{不動点を引く理由}

次の項が
\[
 a_{n+1}=f(a_n),\qquad
 f(z)=\frac{pz+q}{rz+s}
\]
によって決まる場合を考える.
この形の$f$を\textbf{1次分数変換}という.
本節では係数は実数,\,$r\ne0$,\,$ps-qr\ne0$とし,\,元の更新で分母が零にならない範囲を扱う.
$r=0$なら1階の線形の式へ戻り,\,$ps-qr=0$なら定義できる入力で右辺が定数になるため,\,別に調べればよい.

不動点$\alpha$は
\[
 \alpha=\frac{p\alpha+q}{r\alpha+s}
\]
を満たすので,
\[
 r\alpha^2+(s-p)\alpha-q=0
\]
から求まる.
この方程式の根では$r\alpha+s\ne0$である.
もし零なら,\,元の方程式から$p\alpha+q=0$にもなり,\,$ps-qr\ne0$に反する.

不動点の式を使って引き算すると
\[
 f(z)-\alpha
 =\frac{(p-r\alpha)(z-\alpha)}{rz+s}.
\]
定数項は消えるが,\,まだ分母に$z$が残る.
ここで,\,異なるもう一つの不動点$\beta$があれば
\[
 f(z)-\beta
 =\frac{(p-r\beta)(z-\beta)}{rz+s}.
\]
二式の比を取れば,\,共通の分母が消える.
\[
 \frac{f(z)-\alpha}{f(z)-\beta}
 =\frac{p-r\alpha}{p-r\beta}
  \frac{z-\alpha}{z-\beta}.
\]
つまり
\[
 b_n=\frac{a_n-\alpha}{a_n-\beta},\qquad
 \kappa=\frac{p-r\alpha}{p-r\beta}
\]
と置くと,\,$b_{n+1}=\kappa b_n$である.
ここで$p-r\alpha,\,p-r\beta$も非零である.
実際,\,不動点の式から
$(p-r\alpha)(r\alpha+s)=ps-qr$となる.

初項が不動点なら定数列として先に扱う.
$a_1\ne\beta$なら,\,上の差の式から$a_n\ne\beta$が続くため,\,比を取れる.
解いた後は
\[
 b_n=b_1\kappa^{n-1},\qquad
 a_n=\frac{\beta b_n-\alpha}{b_n-1}
\]
で元へ戻す.
$b_n=1$になる場合は有限の$a_n$へ戻せず,\,元の更新もその項まで実数列としては定義できない.

\begin{bookpoint}
一つの不動点を引くと定数項が消える.
二つの不動点からの差の比を取ると,\,残った共通の分母も消える.
置き換えは,\,消したい部分に合わせて組み立てられる.
\end{bookpoint}

\noindent\textbf{【例題】}\par\nobreak
$a_1=3$,\,$a_{n+1}=\frac{2a_n+1}{a_n+2}$を解け.

\noindent\textbf{【方針】}\par\nobreak
不動点1と$-1$からの差の比を取ります.

\noindent\textbf{【解答】}\par\nobreak
\[
 a_n=\frac{2\cdot3^{n-1}+1}{2\cdot3^{n-1}-1}.
\]
\noindent\textbf{【解法】}\par\nobreak
$b_n=\frac{a_n-1}{a_n+1}$とすると
$b_{n+1}=\frac13b_n$,\,$b_1=\frac12$である.
従って$b_n=\frac1{2\cdot3^{n-1}}$となり,
$a_n=\frac{1+b_n}{1-b_n}$から答えを得る.
分母は正で,\,各項も正なので元の分母$a_n+2$は零にならない.
また$b_n\to0$だから$a_n\to1$となる.
不動点からの距離そのものではなく,\,二つの差の比が縮小していた.

\subsection{不動点が重なる場合}

不動点の方程式が重解$\alpha$を持つと,\,異なる二つの差を作れない.
この場合は
\[
 2r\alpha=p-s,\qquad
 p-r\alpha=r\alpha+s=\ell
\]
となり,\,$\ell\ne0$である.
$u_n=a_n-\alpha$と置くと
\[
 u_{n+1}=\frac{\ell u_n}{ru_n+\ell}.
\]
初項が不動点の場合を除き,\,逆数を取って
\[
 \frac1{u_{n+1}}=\frac1{u_n}+\frac r\ell
\]
という等差型へ直せる.
二つの固有方向が分かれない場合に足し算が残った第2章と,\,似た現象が現れた.
その直接の対応は最後の節で確認する.

例えば
\[
 a_1=5,\qquad a_{n+1}=\frac{4a_n-9}{a_n-2}
\]
では$\alpha=3$であり,
\[
 \frac1{a_{n+1}-3}=\frac1{a_n-3}+1.
\]
従って$a_n=3+\frac2{2n-1}$となる.
この例では$a_n>3$が保たれ,\,元の分母も非零である.

\subsection{実数の数列を単位円へ}

不動点の二解が非実数なら,\,実係数なので$\alpha,\overline\alpha$という共役の組になる.
実数$a_n$について
\[
 b_n=\frac{a_n-\alpha}{a_n-\overline\alpha}
\]
と置く.
分子と分母は共役だから,\,$|b_n|=1$である.
また
\[
 \kappa=\frac{p-r\alpha}{p-r\overline\alpha}
\]
も分子と分母が共役なので$|\kappa|=1$となる.
従って,\,$b_{n+1}=\kappa b_n$は\textbf{単位円周上の一定角の回転}である.
実軸上で複雑に見えた動きが,\,円周上では単純になる.

\noindent\textbf{【例題】}\par\nobreak
\[
 a_1=\frac12,\qquad a_{n+1}=\frac{a_n+1}{1-a_n}
\]
を複素平面で解釈し,\,周期を求めよ.

\noindent\textbf{【方針】}\par\nobreak
不動点は$i,-i$です. 差の比を取り,\,一回の回転角を調べます.

\noindent\textbf{【解答】}\par\nobreak
最小周期は4であり,
\[
 \frac12,\ 3,\ -2,\ -\frac13,\ \frac12,\ldots
\]
と繰り返す.

\noindent\textbf{【解法】}\par\nobreak
$b_n=\frac{a_n-i}{a_n+i}$とすると
\[
 b_{n+1}=\frac{1+i}{1-i}b_n=ib_n.
\]
毎回反時計回りに$\frac{\pi}{2}$回転する.
$b_n$は零でなく,\,4回より少ない正の回数では元へ戻らない.
復元は$a_n=i\frac{b_n+1}{1-b_n}$である.
実際に元の式で四項を計算すると上の並びを得る.
どの項も1ではないため,\,すべての更新が定義できる.

ここで,\,回転が定義できることと,\,元の実数列が全項で定義できることは区別する.
例えば同じ式で$a_1=0$とすると$a_2=1$となり,\,$a_3$は定義できない.
円周上では回転を続けられるが,\,復元式の分母が零になる点に到達している.

\subsection{実軸と円を結ぶ座標}\label{節：実軸と円}

回転を三角関数と結び付けるため,\,別の一定倍をした座標
\[
 w=\Phi(a)=\frac{1+ia}{1-ia}
\]
を使おう.
実数$a$なら$|w|=1$であり,\,この$\Phi(a)$は先ほどの$b$に対して$w=-b$である.
例えば$a=0$は$w=1$,\,$a=1$は$w=i$へ移る.

$a=\tan\theta$と書ける範囲では,
\[
 \begin{aligned}
 \Phi(\tan\theta)
 &=\frac{\cos\theta+i\sin\theta}{\cos\theta-i\sin\theta}\\
 &=\cos2\theta+i\sin2\theta.
 \end{aligned}
\]
実数$a$を傾きとして記録した角度が,\,円周上ではその2倍の角度になる.
\NewFigCayley

$w=u+iv$が単位円上で$w\ne-1$なら,\,逆に
\[
 a=\frac{v}{1+u}
\]
で戻れる.
点$(-1,0)$から$(u,v)$への直線の傾きが$a$である.
従って$\Phi$は実軸と,\,一点$-1$を除いた単位円を一対一に対応させる.
除いた一点は,\,傾きが無限に大きくなる方向に当たる.
必要に応じて実軸に無限遠点$\infty$を一つ加えると円全体と対応するが,\,実数の割り算では依然として分母の条件が必要である.

実数$t$を固定し,
\[
 f(a)=\frac{a+t}{1-ta}
\]
とすると,\,直接の計算から
\[
 \Phi(f(a))=\Phi(t)\Phi(a)
\]
となる. $t=\tan\phi$なら$\Phi(t)=e^{2i\phi}$であり,
元の式は$a=\tan\theta$の角度を$\theta+\phi$へ進める式である.
正接の加法公式
\[
 \tan(\theta+\phi)
 =\frac{\tan\theta+\tan\phi}{1-\tan\theta\tan\phi}
\]
が,\,円周上の掛け算として現れた.
角度は$\pi$の整数倍だけ違っても同じ傾きを表す.

\section{特性根が非実数の3項間漸化式}\label{節：非実数の特性根}
\subsection{複素数の冪から実数列へ}

第2章では
\[
 a_{n+2}=s a_{n+1}+t a_n
\]
を,\,特性方程式$X^2-sX-t=0$で調べた.
係数と初期二項が実数でも,\,特性根が非実数になる場合がある.
この場合も二根は異なり,$\lambda,\overline\lambda$という共役の組になる.

虚部が正の根を
\[
 \lambda=\rho e^{i\theta},
 \qquad \rho>0,\quad 0<\theta<\pi
\]
と書く.
二つの冪を用いれば
\[
 a_n=A\lambda^{n-1}+B\overline\lambda^{\,n-1}
\]
である.
根が異なるので,\,初めの二項から$A,B$は一意に決まる.
元の各項は実数だから,\,共役を取った式も同じ数列を表す.
係数の一意性から$B=\overline A$となる.

実数$C,D$を使って
\[
 A=\frac{C-iD}{2},\qquad B=\frac{C+iD}{2}
\]
と書き,$u=(n-1)\theta$と置く.
二つの共役な冪の虚部が打ち消し合い,
\[
 \begin{aligned}
 a_n
 &=\rho^{n-1}\{(A+B)\cos u+i(A-B)\sin u\}\\
 &=\rho^{n-1}(C\cos u+D\sin u)
 \end{aligned}
\]
となる.
従って
\[
 \begin{aligned}
 a_n=\rho^{n-1}\{&
 C\cos((n-1)\theta)\\
 &+D\sin((n-1)\theta)\}.
 \end{aligned}
\]
初期条件から
\[
 C=a_1,\qquad
 D=\cfrac{\cfrac{a_2}{\rho}-a_1\cos\theta}{\sin\theta}.
\]
$0<\theta<\pi$なので$\sin\theta\ne0$であり,\,必ず決められる.

三角関数は,\,実数性を保つために組み合わせた二つの複素数の冪を表している.
「複素数を使ったら最後に偶然実数になる」のではなく,\,共役の対称性によって実数になるのである.

\noindent\textbf{【例題】}\par\nobreak
\[
 a_1=0,\quad a_2=1,\qquad
 a_{n+2}=2a_{n+1}-2a_n
\]
の一般項を求めよ.

\noindent\textbf{【方針】}\par\nobreak
特性根$1\pm i$を極形式で表します.

\noindent\textbf{【解答】}\par\nobreak
\[
 a_n=(\sqrt2)^{n-1}\sin\frac{(n-1)\pi}{4}.
\]
\noindent\textbf{【解法】}\par\nobreak
特性方程式は$X^2-2X+2=0$である.
$1+i=\sqrt2e^{i\frac{\pi}{4}}$なので,
$\rho=\sqrt2,\,\theta=\frac{\pi}{4}$となる.
$a_1=0$から$C=0$,
$a_2=1$から$1=\sqrt2D\sin\frac{\pi}{4}=D$と決まる.
最初の数項は$0,1,2,2,0,-4,\ldots$であり,\,得た式とも一致する.
証明は数項の一致ではなく,\,初期二項と更新の一致による.

\subsection{回転を直接記録する複素数}\label{節：複素数の状態}

前の計算では解を冪に分けた.
今度は,\,元の二項から回転を表す複素数を直接作ろう.
根と係数の関係から
\[
 s=\lambda+\overline\lambda,\qquad
 t=-\lambda\overline\lambda.
\]
そこで
\[
 z_n=a_{n+1}-\overline\lambda a_n
\]
と置けば,
\[
 \begin{aligned}
 z_{n+1}
 &=a_{n+2}-\overline\lambda a_{n+1}\\
 &=\lambda a_{n+1}-\lambda\overline\lambda a_n\\
 &=\lambda z_n.
 \end{aligned}
\]
実数列から作った$z_n$は,\,毎回$\rho$倍されながら$\theta$回転する.
第2章で作った補助数列$a_{n+1}-\lambda_1a_n$を,\,複素数の範囲まで広げたのである.

$\lambda=u+iv$で$v>0$と書くと,
\[
 z_n=(a_{n+1}-ua_n)+iv a_n.
\]
従って
\[
 a_n=\frac{\operatorname{Im}z_n}{v},\qquad
 a_{n+1}=\operatorname{Re}z_n+u a_n.
\]
二つの実数と一つの複素数は一対一に対応し,\,情報は失われていない.
元の状態$(a_{n+1},a_n)$から座標を変えたら,\,回転と拡大・縮小が見えたのである.

例えば先ほどの式では$u=v=1$だから
\[
 z_n=(a_{n+1}-a_n)+ia_n.
\]
$z_1=1$より$z_n=(1+i)^{n-1}$となり,\,虚部を取れば一般項が得られる.
\NewFigSpiral

\subsection{元の座標では円とは限らない}

$z_n$の平面では,$\rho=1$なら原点からの距離が一定になる.
しかし元の$(a_{n+1},a_n)$の平面では,\,座標変換が一般に長さや直角を保たないため,\,円上にあるとは限らない.
対応する曲線は
\[
 (a_{n+1}-ua_n)^2+v^2a_n^2=\text{一定}
\]
という楕円になる.
更新が通常の座標でそのまま回転行列である,\,という意味ではない.
実際の状態はその楕円上の点の一部をたどる.

$\rho\ne1$では上の量が毎回$\rho^2$倍される.
この量は,\,元の座標で見た回転の半径に相当する.
固有値が非実数のときに実固有方向が見つからなかったのは,\,固定した一本の実方向にとどまらず,\,平面を回っているからである.

\subsection{周期と収束}

$0<\rho<1$なら,\,正弦と余弦は有界なので$a_n\to0$となる.
$\rho=1$なら有界な振動となるが,\,必ず周期的とは限らない.
非零の実数解が周期$T$を持つための条件は
\[
 \rho=1,\qquad T\theta\in2\pi\mathbb Z
\]
である. ここで$\mathbb Z$は整数全体を表す.

十分性は$\lambda^T=1$から分かる.
必要性は,\,周期があると状態から作った$z_n$にも$z_{n+T}=z_n$が成り立つことから分かる.
非零解では$z_1\ne0$なので$\lambda^Tz_1=z_1$から$\lambda^T=1$となる.
従って周期を持つのは,$\rho=1$かつ$\frac{\theta}{2\pi}$が有理数のときである.
零数列は別に扱う.

$\rho>1$なら非零解の状態は拡大する.
ただし$a_n$だけには零になる項もあり,\,各項の絶対値が単調に増えるとは限らない.
先ほどの例でも4項ごとに零が現れる.
\textbf{状態の大きさと,\,その一成分の大きさを区別する}.

\section{三角関数と漸化式}\label{節：三角関数と漸化式}
\subsection{一定角を足す場合}

単位円上の点を$z_n=e^{i\theta_n}$と書き,\,一定角$\phi$ずつ進めると
\[
 z_{n+1}=e^{i\phi}z_n,\qquad
 \theta_n=\theta_1+(n-1)\phi
\]
となる.
横座標$a_n=\cos\theta_n$だけを記録すると,\,加法公式から
\[
 \cos(\theta_n+2\phi)+\cos\theta_n
 =2\cos\phi\cos(\theta_n+\phi)
\]
なので
\[
 a_{n+2}=2\cos\phi\,a_{n+1}-a_n.
\]
一つの横座標だけでは点の上下を区別できないが,\,隣り合う二つの記録があれば次を更新できる.
特性方程式の根は$e^{i\phi},e^{-i\phi}$であり,\,前節につながる.
$0<\phi<\pi$では二根が非実数となる.

一方,\,傾き$a_n=\tan\theta_n$を記録すれば
\[
 a_{n+1}=\frac{a_n+\tan\phi}{1-a_n\tan\phi}
\]
という1次分数型になる.
ただし$\tan\phi$と各項の正接が定義され,\,更新の分母が非零の範囲で考える.
同じ角度の動きでも,\,\textbf{横座標と傾きでは更新式の形が変わる}.

\subsection{角度を倍にする場合}

次は
\[
 a_1=\frac{\sqrt3}{2},\qquad a_{n+1}=2a_n^2-1
\]
を考える.
倍角公式$\cos2\theta=2\cos^2\theta-1$と,
$a_1=\cos\frac{\pi}{6}$を使えば,
\[
 a_n=\cos\left(2^{n-1}\frac{\pi}{6}\right)
\]
と予想できる.
初項に合い,\,この式を更新の右辺に代入すると次の項の式になるので,\,帰納法で正しい.

角度は$\theta_{n+1}=2\theta_n$と選べる.
複素数では
\[
 z_{n+1}=z_n^2,\qquad z_n=z_1^{\,2^{n-1}}
\]
であり,$a_n=\operatorname{Re}z_n$である.
\NewFigDoubling

これは一定角の回転ではない.
現在の角度$\theta_n$をもう一度足すため,\,毎回の回転角が変わる.
さらに$z$と$-z$は同じ$z^2$へ移るので,\,円周全体の写像として一対一ではない.

余弦の値から角度は一意に決まらない.
そのため$a_{n+1}=2a_n^2-1$から,\,選び方によらず$\theta_{n+1}=2\theta_n$が強制されると読むのは正確でない.
整合する角度の列を一つ選び,\,そこから作った候補が元の式を満たすことを示している.

\subsection{正接の倍角と定義域}

傾きを記録する倍角型は
\[
 a_{n+1}=\frac{2a_n}{1-a_n^2}
\]
である.
$a_n=\tan\theta_n$と表せれば,$\theta_{n+1}=2\theta_n$と選べる.
ただし$a_n=\pm1$では次の傾きが縦向きになり,\,実数の値としては定義できない.
円周上の動きが続いても,\,傾きの列が続くとは限らない.

\noindent\textbf{【例題】}\par\nobreak
\[
 a_1=2-\sqrt3,\qquad
 a_{n+1}=\frac{2a_n}{1-a_n^2}
\]
の一般項を求め,\,全項が定義されることを示せ.

\noindent\textbf{【方針】}\par\nobreak
$2-\sqrt3=\tan\frac{\pi}{12}$です. 角度を倍にし,\,途中で縦向きにならないか確認します.

\noindent\textbf{【解答】}\par\nobreak
\[
 a_n=\tan\left(2^{n-1}\frac{\pi}{12}\right).
\]
\noindent\textbf{【解法】}\par\nobreak
初項と倍角公式から候補を得る.
$n=1$では正接は定義できる.
$n\ge2$でその角度が$\frac{\pi}{2}+\ell\pi$に一致するとすれば
\[
 2^{n-2}=3+6\ell
\]
となるが,2の累乗は3の倍数でないので不可能である.
従って全項の正接が定義される.
また次の角度も縦向きにならないため,\,更新の分母$1-a_n^2$も非零である.
帰納法により一般項が成立する.

なお,$\Phi$は前節の実軸と円の対応であり,
\[
 \Phi\left(\frac{2a}{1-a^2}\right)=\Phi(a)^2
\]
となる.
正接の倍角型も,\,円周上の二乗を別の座標で記録していた.
1次分数型の一定角回転とは,\,同じ単位円を使っても更新の規則が違う.

\subsection{整数倍角とチェビシェフ多項式}

角度を2倍する式を見たので,\,一般の整数倍も考えよう.
多項式$T_m(x)$を
\[
 T_0(x)=1,\quad T_1(x)=x,
\]
\[
 T_{m+1}(x)=2xT_m(x)-T_{m-1}(x)
\]
で定める. これを\textbf{チェビシェフ多項式}という.
加法公式によって
\[
 T_m(\cos\theta)=\cos(m\theta)
\]
が帰納的に示される.
例えば
\[
 T_2(x)=2x^2-1,\qquad T_3(x)=4x^3-3x.
\]
特定の$x$を固定すると,$m$に関しては2階の線形漸化式である.

一方,\,固定した整数$d\ge2$について
\[
 a_{n+1}=T_d(a_n)
\]
を反復すると,$a_1=\cos\theta$のとき
\[
 a_n=\cos(d^{n-1}\theta)
\]
となる. こちらは$n$に関して非線形の更新である.
多項式の次数を進める添字$m$と,\,更新回数を表す$n$を区別しよう.

さらに
\[
 T_j(T_m(x))=T_{jm}(x)
\]
という合成の公式もある.
$x=\cos\theta$で確かめると両辺は$\cos(jm\theta)$になる.
両辺は多項式で,\,区間$[-1,1]$の無数の点で一致するため,\,すべての実数・複素数$x$で一致する.
零でない多項式の根が次数を超えないことを使っている.
従って一般の初項でも
\[
 a_n=T_{d^{n-1}}(a_1)
\]
と表せる.

\subsection{三角関数の外側にも続く式}

実数$a_1$が$[-1,1]$の外にあると,\,実角の余弦としては表せない.
しかし多項式の更新自体は定義されている.
任意の$x$に対し,\,非零の複素数$u$を
\[
 u^2-2xu+1=0
\]
の根として選べば,$x=\frac{u+u^{-1}}2$であり,
\[
 T_m(x)=\frac{u^m+u^{-m}}2
\]
となる. 初めの二つの値と更新を確かめればよい.
$|u|=1$なら余弦の表示に戻り,\,実数$x>1$なら実根を選んで指数的な増大として読める.

例えば$a_1=\frac54$,$a_{n+1}=2a_n^2-1$では$u=2$と選べるので,
\[
 a_n=\frac{2^{\,2^{n-1}}+2^{-2^{n-1}}}{2}.
\]
三角関数の公式は,\,二つの逆数の冪を足した式の一つの姿でもあった.
第2章の等比的な成分への分解と,\,ここでも再会する.

\section{回転から見直す漸化式}\label{節：回転から見直す}
\subsection{複素数の掛け算と行列}

複素数$u+iv$を$x+iy$に掛けると
\[
 (u+iv)(x+iy)=(ux-vy)+i(vx+uy).
\]
実部と虚部を縦に並べれば,
\[
 \begin{pmatrix}x'\\y'\end{pmatrix}
 =
 \begin{pmatrix}u&-v\\v&u\end{pmatrix}
 \begin{pmatrix}x\\y\end{pmatrix}.
\]
$u=\rho\cos\theta,\,v=\rho\sin\theta$なら,\,これは
\[
 \rho
 \begin{pmatrix}
 \cos\theta&-\sin\theta\\
 \sin\theta&\cos\theta
 \end{pmatrix}
\]
である. 複素数の掛け算と,\,二成分の実数の行列更新が一致した.

この行列の特性方程式は
\[
 X^2-2uX+(u^2+v^2)=0
\]
で,\,根は$u\pm iv$である.
実数の行列に非実数の固有値が現れることは,\,回転の倍率を複素数で記録したことに対応する.

ただし非実数の固有値を持つ一般の実行列が,\,元の座標でこの形に見えるとは限らない.
\sectionref{節：複素数の状態}のように,\,座標を変えて初めて回転と拡大・縮小になる場合がある.

\subsection{1次分数型も行列から生まれる}\label{節：分数と行列の合流}

1次分数型
\[
 a_{n+1}=\frac{pa_n+q}{ra_n+s}
\]
を,\,今度は二つの数の比で持ってみよう.
$(x_1,y_1)=(a_1,1)$から
\[
 \begin{pmatrix}x_{n+1}\\y_{n+1}\end{pmatrix}
 =M\begin{pmatrix}x_n\\y_n\end{pmatrix},
 \qquad
 M=\begin{pmatrix}p&q\\r&s\end{pmatrix}
\]
と更新すれば,
\[
 \frac{x_{n+1}}{y_{n+1}}
 =\cfrac{p\frac{x_n}{y_n}+q}{r\frac{x_n}{y_n}+s}.
\]
従って,\,分母が非零の範囲では$a_n=\frac{x_n}{y_n}$となる.
一次の連立更新を比として記録すると,\,1次分数型が生まれるのである.

二つの成分を同じ非零数で倍しても比は変わらない.
長さと矢印の正負を区別せず,\,原点を通る直線としての方向を記録するこの考え方を,\,\textbf{射影的な見方}という.
分母$y_n=0$は,\,成分比を有限の実数で記録できなくなる方向に当たる.
元の実数列で定義できる範囲と,\,行列で状態を追える範囲は区別する.

不動点$\alpha$については
\[
 M\begin{pmatrix}\alpha\\1\end{pmatrix}
 =(r\alpha+s)\begin{pmatrix}\alpha\\1\end{pmatrix}.
\]
従って,\,\textbf{不動点は固有値ではなく,\,固有ベクトルの成分比}である.
対応する固有値は$\lambda_\alpha=r\alpha+s$であり,\,役割が異なる.
\NewFigProjective

\subsection{差の比は固有成分の比}

異なる二つの不動点$\alpha,\beta$に対し,\,状態を
\[
 \begin{pmatrix}x_n\\y_n\end{pmatrix}
 =c_n\begin{pmatrix}\alpha\\1\end{pmatrix}
 +d_n\begin{pmatrix}\beta\\1\end{pmatrix}
\]
と分ける. 固有値を$\lambda_\alpha,\lambda_\beta$とすれば
\[
 c_{n+1}=\lambda_\alpha c_n,\qquad
 d_{n+1}=\lambda_\beta d_n.
\]
成分を比較すると
\[
 x_n-\alpha y_n=(\beta-\alpha)d_n,
\]
\[
 x_n-\beta y_n=(\alpha-\beta)c_n.
\]
従って,\,比が定義できる場合には
\[
 \frac{a_n-\alpha}{a_n-\beta}=-\frac{d_n}{c_n}
\]
となり,\,その公比は$\frac{\lambda_\beta}{\lambda_\alpha}$である.
不動点の和$r(\alpha+\beta)=p-s$から
\[
 \lambda_\beta=p-r\alpha,\qquad
 \lambda_\alpha=p-r\beta
\]
なので,\,前の1次分数型の公比と一致する.

本節でも$r\ne0,\,ps-qr\ne0$を仮定する.
このとき$M$は逆行列を持ち,\,固有値も非零である.
不動点を初項とする定数列を除けば,\,必要な成分比を取り続けられる.
元の更新の分母が零にならない条件は別に確かめる.

実係数で不動点が共役な非実数なら,\,固有値も共役である.
両者は絶対値が等しいため,\,\textbf{比を取ると共通の拡大率が消えて回転だけが残る}.
3項間漸化式では回転に$\rho^{n-1}$が伴ったが,\,1次分数型の円周上の座標ではその大きさを除いていた.

不動点が重なる場合は,\,行列も二つの固有方向に分けられない.
$r\ne0$だから$M$は単位行列の定数倍ではなく,\,固有方向は一本になる.
その寄与の蓄積を比で読み直すと,\,逆数を取った等差型が現れる.
第2章の重解と,\,第3章の重なった不動点がここでつながった.

\subsection{一つの例を三通りに記録する}

特に
\[
 M=\begin{pmatrix}1&1\\-1&1\end{pmatrix}
\]
を使おう.
状態の更新は$x_{n+1}=x_n+y_n,\,y_{n+1}=-x_n+y_n$である.
複素数を$z_n=y_n+ix_n$と選ぶと,
\[
 z_{n+1}=(1+i)z_n.
\]
この座標では,\,毎回$\sqrt2$倍されながら反時計回りに$\frac{\pi}{4}$回転する.
各成分は特性方程式$X^2-2X+2=0$に対応し,
\[
 x_{n+2}=2x_{n+1}-2x_n
\]
などを満たす.

一方,\,成分比$a_n=\frac{x_n}{y_n}$を記録すると,
\[
 a_{n+1}=\frac{a_n+1}{1-a_n}.
\]
共通の拡大率が消え,\,傾きの変化だけが残る.
さらに$w_n=\Phi(a_n)$で円周へ移すと,
\[
 w_{n+1}=iw_n.
\]
角度を2倍して記録する$\Phi$によって,\,回転角は$\frac{\pi}{2}$になる.

$(x_1,y_1)=(\frac12,1)$なら,\,成分比は先ほどの4周期の列となる.
行列では4回後の状態は$M^4\boldsymbol v_1=-4\boldsymbol v_1$であり,\,元の状態そのものには戻らない.
しかし両成分が同じ倍率を受けるため比は戻る.
\textbf{状態の周期と,\,比の周期は同じとは限らない}.

\subsection{何を記録しているか}

同じ動きを別の形で記録できることを確かめた.
一方で,\,式の見た目が似ていても更新の仕組みが異なる場合もある.

\begin{itemize}
 \item $z_{n+1}=\lambda z_n$は,\,一定角の回転と一定倍率の拡大・縮小である.
 \item $z_{n+1}=z_n^d$は,\,偏角を$d$倍し,\,絶対値も$d$乗する. 単位円上でも一定角の回転とは異なる.
 \item 実部や傾きだけを記録すると,\,異なる点が同じ値になることがある. 何の情報を残し,\,何を区別しなくなったかを確認する.
\end{itemize}

複素平面への置き換えが,\,常に全体の座標変換になるわけではない.
例えば余弦は上下二点を同じ横座標に送る.
それでも,\,元の更新と整合する点の動きを一つ作れば,\,数列の式を導ける.
一対一の座標変換と,\,一部の情報を読み取る操作を区別して使おう.

\begin{bookpoint}
漸化式の形は,\,動きそのものだけでなく,\,\textbf{その動きの何を数として記録したか}によって決まる.
不動点からのずれ,\,固有方向の成分,\,成分比,\,実部,\,傾き.
追う量を選び直すと,\,別々に見えた解法のつながりが見えてくる.
\end{bookpoint}

\end{multicols}
